\documentclass[12pt,letterpaper]{article}

\usepackage[utf8]{inputenc}
\usepackage[T1]{fontenc}
\usepackage[spanish]{babel}
\usepackage{mathptmx} 
\usepackage{graphicx}
\usepackage{float}          %Para poder insertar imágenes o tablas de manera flotante (y que aparezcan en donde uno desea) 
%\geometry{top=2.5cm, bottom=2.5cm, left=2.5cm, right=2.5cm} %Márgenes
\usepackage{amsmath,amssymb}
\usepackage{geometry}
\geometry{margin=2.5cm}

\begin{document}


\hrule
\vspace{1mm}
\hrule
\vspace{4mm}

\begin{center}
     \Huge \textbf{CONTROL 1}\\[0.3cm]
     \Large GEOMETRÍA EN EL PLANO Y ESPACIO\\[0.3cm]
     \Large PEDAGOGÍA EN MATEMÁTICA Y COMPUTACIÓN\\[0.2cm]
     \large UNIVERSIDAD DE SANTIAGO DE CHILE\\[0.3cm]
    
     \normalsize Angela Meza\\
     \normalsize 27 de Octubre de 2026\\
\end{center}

\vspace{5mm}
\hrule
\vspace{1mm}
\hrule
\vspace{1mm}

\section*{Introducción}
\begin{itemize}
    \item Redactar una demostración para la proposición:  Sea el ángulo $\angle{AOB}$, entonces $l$ es bisectriz de $\angle{AOB}$
    $l$ equidista de los lados de $\angle{AOB}$
\end{itemize}

\section*{PARTE 1: Caracterización de la Bisectriz}

Dado el ángulo $\angle AOB$, consideramos la recta $l$ como su bisectriz, por lo que se cumple que $\angle(OA, l) = \angle(l, OB)$. 

Nuestro objetivo es demostrar que la recta $l$ es equidistante a los lados del ángulo.

\begin{figure}[H]
    \centering
    \includegraphics[width=0.5\linewidth]{imagen 1.png}
    \caption{Construcción de la bisectriz y proyecciones de $P$ sobre los lados del ángulo.}
    \label{fig:placeholder}
\end{figure}

\end{document}